\phantomsection
\chapter{Проектування та експериментальні дослідження}
\label{chapter:development-and-analysis}

У цьому розділі наведено практичні приклади оформлення звітної документації згідно з вимогами ДСТУ~3008:2015. Розглянуто ієрархію заголовків різних рівнів, варіанти масштабування та позиціонування графічних матеріалів, оформлення лістингів вихідного коду, а також створення багатосторінкових таблиць із автоматичним переносом шапки.

\section{Ієрархія заголовків та робота з графікою}
\label{sec:headings-and-graphics}

Текст кожного підрозділу повинен починатися з абзацного відступу розміром 5 символів (5ex). Згідно з вимогами стандарту, між заголовками однакового або суміжних рівнів витримується нормативний інтервал без додаткових порожніх рядків.

\subsection{Варіанти масштабування зображень}
\label{subsec:image-scaling}

Всі рисунки мають бути вирівняні по центру сторінки та супроводжуватися підрисунковим підписом виду «Рисунок X.Y~--~Назва рисунка». Для вставки графіки рекомендується використовувати макрос \texttt{\textbackslash paddedimage}, який забезпечує нормативні відступи зверху та знизу.

Нижче продемонстровано великомасштабне зображення (масштаб 0.85), яке підходить для деталізованих блок-схем, діаграм класів або загальних планів архітектури.

\begin{figure}[H]
    \centering
    \paddedimage[width=0.85\textwidth]{cat-1-wide.png}
    \caption{Загальна архітектура інформаційної взаємодії системи (великий масштаб)}
    \label{fig:arch-wide}
\end{figure}

Для компактних схем, графіків функцій чи окремих елементів інтерфейсу доцільно обирати середній масштаб (0.52), як показано на рис.~\ref{fig:component-medium}.

\begin{figure}[H]
    \centering
    \paddedimage[width=0.52\textwidth]{cat-2-square.jpg}
    \caption{Структурна схема модуля обробки запитів (середній масштаб)}
    \label{fig:component-medium}
\end{figure}

У випадках, коли необхідно проілюструвати допоміжний графічний елемент, піктограму чи компактний стан системи, використовується малий масштаб (0.30).

\begin{figure}[H]
    \centering
    \paddedimage[width=0.30\textwidth]{cat-2-square.jpg}
    \caption{Піктограма статусу обчислювального потоку (малий масштаб)}
    \label{fig:icon-small}
\end{figure}

\subsection{Способи посилання на графічні матеріали}
\label{subsec:image-refs}

У звітах згідно з ДСТУ допускається кілька варіантів посилань на ілюстрації в тексті:
\begin{dashlist}[wide,beginpenalty=10000]
    \item Пряме посилання через команду \texttt{\textbackslash ref}: «як наведено на рис.~\ref{fig:arch-wide}».
    \item Автоматичне посилання через пакет \texttt{cleveref} (\texttt{\textbackslash cref}): «згідно з \cref{fig:component-medium}».
    \item Посилання на початку речення з великої літери (\texttt{\textbackslash Cref}): «\Cref{fig:arch-wide} демонструє ключові компоненти».
    \item Посилання в дужках: «схема ініціалізації середовища (див. рис.~\ref{fig:icon-small})».
    \item Посилання на діапазон рисунків (\texttt{\textbackslash crefrange}): «динаміка змін відображена на \crefrange{fig:arch-wide}{fig:icon-small}».
\end{dashlist}

\section{Вставка та оформлення вихідного коду}
\label{sec:code-listings}

Лістинги програмного коду оформлюються за допомогою пакету \texttt{listings}. Заголовок лістингу розташовується знизу або зверху відповідно до вимог кафедри, кожна строка нумерується, а ключові слова мови програмування виділяються напівжирним шрифтом або кольором.

\begin{lstlisting}[label={lst:quicksort}]
def quick_sort(arr: list) -> list:
    """Рекурсивна реалізація алгоритму швидкого сортування."""
    if len(arr) <= 1:
        return arr
    pivot = arr[len(arr) // 2]
    left = [x for x in arr if x < pivot]
    middle = [x for x in arr if x == pivot]
    right = [x for x in arr if x > pivot]
    return quick_sort(left) + middle + quick_sort(right)
\end{lstlisting}

Посилання на лістинг у тексті здійснюється стандартним чином: «у лістингу~\ref{lst:quicksort} наведено реалізацію алгоритму» або «відповідно до \cref{lst:quicksort} базовим випадком рекурсії є одноелементний список».

\section{Багатосторінкові таблиці}
\label{sec:tables}

\subsection{Оформлення таблиць згідно з ДСТУ}
\label{subsec:tables-dstu}

Таблиці у шаблоні реалізовано на базі середовища \texttt{longtable}. Команда створення таблиці приймає чотири обов'язкові аргументи:
\begin{enumerate}[wide,beginpenalty=10000]
    \item Специфікація вирівнювання колонок (наприклад, \texttt{\{|c|p\{7.5cm\}|c|c|\}}).
    \item Назва таблиці (наприклад, \texttt{Результати експерименту}).
    \item Мітка для посилань (наприклад, \texttt{tab:benchmarks}).
    \item Шапка таблиці (перший рядок заголовків стовпців із горизонтальними лініями).
\end{enumerate}

\subsubsection{Приклад таблиці, що переноситься на наступну сторінку}
\label{subsubsec:multipage-table}

Якщо таблиця не вміщується на одній сторінці, шаблон автоматично розриває її та додає на наступній сторінці напис «Продовження таблиці X.Y» разом із повторенням шапки таблиці, що повністю відповідає нормам ДСТУ~3008:2015.

Зведені показники порівняльного аналізу алгоритмів наведено в табл.~\ref{tab:benchmarks}. Зверніть увагу: \cref{tab:benchmarks} займає більше однієї сторінки, демонструючи автоматичний механізм розриву.

\begin{table}{|c|p{7.5cm}|c|c|}
{Результати комплексного порівняльного аналізу параметрів системи}
{tab:benchmarks}
{\hline № & Найменування тестового параметра & Метод А & Метод Б \\ \hline}
1 & Час первинної ініціалізації обчислювального середовища (мс) & 14.2 & 9.8 \\ \hline
2 & Базовий обсяг оперативної пам'яті для кешування (МБ) & 64.0 & 48.5 \\ \hline
3 & Пропускна здатність вхідного потоку транзакцій (тис./с) & 1250 & 1480 \\ \hline
4 & Затримка обробки критичних повідомлень під час піку (мс) & 3.2 & 2.1 \\ \hline
5 & Рівень навантаження на обчислювальні ядра процесора (\%) & 78.4 & 62.1 \\ \hline
6 & Середнє питоме енергоспоживання вузла кластера (Вт) & 46.2 & 38.0 \\ \hline
7 & Часова складність обробки у найкращому сценарії & $O(n)$ & $O(1)$ \\ \hline
8 & Часова складність обробки в середньому сценарії & $O(n \log n)$ & $O(n)$ \\ \hline
9 & Часова складність обробки у найгіршому сценарії & $O(n^2)$ & $O(n \log n)$ \\ \hline
10 & Просторова складність алгоритмічної моделі & $O(n)$ & $O(1)$ \\ \hline
11 & Стабільність збереження порядку пар еквівалентних ключів & Так & Ні \\ \hline
12 & Рівень горизонтального масштабування обчислень & Високий & Помірний \\ \hline
13 & Чутливість структури даних до фрагментації пам'яті & Низька & Середня \\ \hline
14 & Частота влучань у дворівневий кеш процесора L1/L2 (\%) & 94.1 & 89.6 \\ \hline
15 & Швидкість деградації продуктивності при апаратних збоях & 3.8\% & 6.5\% \\ \hline
16 & Кількість операцій порівняння на контрольний блок (тис.) & 8.45 & 6.20 \\ \hline
17 & Кількість операцій перестановки на блок (тис.) & 2.10 & 1.85 \\ \hline
18 & Час обробки масиву розмірністю $N = 10^4$ (мс) & 4.2 & 3.1 \\ \hline
19 & Час обробки масиву розмірністю $N = 10^5$ (мс) & 48.6 & 35.8 \\ \hline
20 & Час обробки масиву розмірністю $N = 10^6$ (мс) & 580.2 & 418.4 \\ \hline
21 & Час обробки попередньо впорядкованого масиву (мс) & 12.1 & 8.2 \\ \hline
22 & Час обробки масиву, впорядкованого у зворотному напрямку (мс) & 710.4 & 538.1 \\ \hline
23 & Вплив типу даних на загальну тривалість виконання & Низький & Помірний \\ \hline
24 & Надійність збереження стану транзакцій при відмові (\%) & 99.99 & 99.95 \\ \hline
25 & Загальний приріст продуктивності після оптимізації & +34.5\% & +41.8\% \\ \hline
\end{table}

\subsubsection{Аналіз отриманих табличних даних}
\label{subsubsec:table-analysis}

Аналіз даних, наведених у табл.~\ref{tab:benchmarks}, свідчить про те, що Метод Б забезпечує перевагу у швидкодії на великих обсягах вибірки ($N \ge 10^5$), проте Метод А демонструє стабільність порядку елементів та вищу стійкість до фрагментації пам'яті.
